\begin{frame}{이번 주차 정리}
  \begin{enumerate}
    \item \kb{14.1 PCA} --- ``분산을 최대화하라''는 말과 ``가장 큰
      고유벡터를 찾아라''는 말은 \textbf{같은 문제} (라그랑주
      승수법 $\to$ $\Sigma v = \lambda v$). \; 주성분의 분산 =
      \textbf{고유값}, 버린 정보량 = \textbf{버린 고유값의 합}.
      \; 단위 다르다면 \textbf{반드시 표준화} 후 PCA.
    \item \kb{14.2 word2vec} --- 분포 가설(``친구를 보면 그 단어를
      안다'')을 \textbf{이진 분류}(네거티브 샘플링)로 학습 ---
      softmax $O(V)$를 $O(k_{\text{neg}})$로 대체. \; 학습된 벡터에서
      \textbf{왕 $-$ 남자 $+$ 여자 $\approx$ 여왕}: 의미 축의
      직교 + 스케일 일관으로 닫히는 병렬 4각형.
    \item \kb{14.3 임베딩 실전} --- Transformer의 첫 층도 \textbf{학습된
      임베딩 테이블}, RAG = \textbf{질문도 벡터, 문서도 벡터, 코사인
      유사도로 가까운 것을 검색} --- 14.2의 단어 유사도와 \textbf{같은
      연산}.
  \end{enumerate}
  \vskip0.8em
  \begin{block}{다음 주 --- Week 15. 잠재변수 생성모델}
    이 장은 관측 데이터 속에 숨은 저차원 구조를 \textbf{찾는} 쪽 ---
    15장은 그 구조(잠재 공간)가 데이터를 어떻게 \textbf{만들어내는가}
    를 학습하는 생성 모델: \kb{VAE} $\to$ \kb{GAN}·\kb{Diffusion}.
  \end{block}
\end{frame}
